% ---
% Reemplaza líneas 1513-1584 de Docs/Technical_Report.tex
% Decisiones: N/A (sin cambio)
% ---

\subsection{Red neuronal clásica e integración
            Qiskit--PyTorch}

% -----------------------------------------------------------
\subsubsection{Perceptrón multicapa (MLP)}
% -----------------------------------------------------------

La red neuronal clásica utilizada como clasificador
\textit{downstream} en el pipeline del presente trabajo es
un perceptrón multicapa (MLP, \textit{Multi-Layer
Perceptron}), una arquitectura de aprendizaje supervisado
compuesta por capas de neuronas totalmente conectadas
\cite{goodfellow2016deep}.

Formalmente, dado un vector de entrada
$\mathbf{z} \in \mathbb{R}^n$, que puede ser el
\textit{quantum embedding}
$\mathbf{e}(\mathbf{x}, \boldsymbol{\theta})$ o cualquiera
de los vectores de las condiciones clásicas, la
transformación de la capa $\ell$ se define como:

\begin{equation}
  \mathbf{h}^{(\ell)} = \sigma\!\left(
  W^{(\ell)}\,\mathbf{h}^{(\ell-1)} + \mathbf{b}^{(\ell)}
  \right), \qquad \ell = 1, \ldots, L
  \label{eq:mlp_layer}
\end{equation}

\noindent donde $W^{(\ell)} \in \mathbb{R}^{d_\ell \times
d_{\ell-1}}$ es la matriz de pesos, $\mathbf{b}^{(\ell)}
\in \mathbb{R}^{d_\ell}$ es el vector de sesgos, y
$\sigma(\cdot)$ es la función de activación no lineal.
La capa de salida produce una probabilidad de clasificación
binaria mediante la función sigmoide:

\begin{equation}
  \hat{y} = \sigma_{\text{sig}}\!\left(
  W^{(L)}\,\mathbf{h}^{(L-1)} + b^{(L)}
  \right) =
  \frac{1}{1 + e^{-(W^{(L)}\mathbf{h}^{(L-1)}
  + b^{(L)})}}
  \in [0, 1]
  \label{eq:sigmoid_output}
\end{equation}

\noindent El modelo se entrena minimizando la función de
pérdida de entropía cruzada binaria (\textit{Binary
Cross-Entropy}, BCE):

\begin{equation}
  \mathcal{L}_{\text{BCE}}(\hat{y}, y) =
  -\left[y\log\hat{y} + (1-y)\log(1-\hat{y})\right]
  \label{eq:bce_loss}
\end{equation}

\noindent donde $y \in \{0, 1\}$ es la etiqueta real
(0\,=\,benigno, 1\,=\,maligno). Los parámetros
$\{W^{(\ell)}, \mathbf{b}^{(\ell)}\}$ se optimizan
mediante descenso de gradiente estocástico (\textit{Adam}
o equivalente) a través del algoritmo de
\textit{backpropagation} \cite{goodfellow2016deep}.

La arquitectura del MLP se mantiene idéntica para las cinco
condiciones experimentales: mismo número de capas ocultas,
mismas dimensiones, misma función de activación, mismos
hiperparámetros y misma semilla aleatoria. Esta
invarianza es el principio metodológico central que garantiza
que las diferencias de rendimiento entre condiciones sean
atribuibles únicamente a la transformación aplicada en M5,
y no a diferencias en el clasificador.
